\documentclass{situationist-preprint}

\title{The Justice of Markets}
\msid{POL-12-a}
\msfamily{POL-12}
\msdiscipline{Political Philosophy}
\mstype{Principal manuscript}
\mskeywords{markets; distributive justice; difference principle; basic structure; entitlement; exploitation; Pareto efficiency; noxious markets; background justice}

\begin{document}
\maketitle

\begin{abstract}
Are markets just? Libertarians answer that market outcomes are just whatever their pattern, because they arise from voluntary exchange; socialists answer that markets are unjust in themselves, because they rest on exploitation and reward fortune; others hold that justice does not apply to the impersonal outcomes of markets at all. We argue that markets are just when their outcomes satisfy the principles of justice defended in POL-01, equal basic liberty and the modified difference principle, and that many, but not all, actual markets do. We make the criterion precise by treating a market as a component of the basic structure, defined by its property, contract and regulatory rules, and by assessing it against the alternatives obtainable by changing those rules with the other institutions held fixed. We prove that the justice of the whole regime implies the justice of each market in it but not conversely, that justice implies Pareto efficiency but not conversely, and that the voluntariness of exchanges is neither necessary nor sufficient for justice. We then apply the criterion to Nozick's Wilt Chamberlain case, sweatshop labour, markets in votes and organs, famines and winner-take-all labour markets, and answer the main libertarian, Hayekian, socialist and expressivist objections.
\end{abstract}

\section{The question}

Markets allocate most of the goods and much of the labour of modern societies. They decide, through the prices at which people are willing to buy and sell, who produces what, who works where and who earns how much. The question whether markets are just is therefore not a question about an isolated institution. It is a large part of the question whether the societies in which we live are just.

Three answers dominate the literature. The first holds that market outcomes are just whatever their pattern, provided that they arise by voluntary exchange from a just starting point. On \citeauthor{Nozick1974}'s entitlement theory \citeyearpar{Nozick1974}, whatever arises from a just situation by just steps is itself just, and voluntary exchange is a just step; any principle that requires a particular pattern of distribution will be upset by free exchange and can be maintained only by continual interference with liberty. The second holds that markets are unjust in themselves. On the Marxian view, the wage relation is exploitative, since it allows the owners of capital to appropriate the surplus produced by labour \citep{Marx1867}; on \citeauthor{Cohen2009}'s view \citeyearpar{Cohen2009}, markets motivate production by greed and fear and distribute by a principle that rewards luck, and a just and communal society would replace them where it could. The third holds that justice does not apply to market outcomes at all. \citet{Hayek1976} argued that \enquote{social justice} is a category mistake when applied to the results of a spontaneous order, since injustice requires an agent who intended the result, and no one designs the distribution that a market produces.

We defend a fourth answer, which has its origin in Rawls. Markets are just when their outcomes satisfy the principles of justice of POL-01: a principle of equal basic liberties and the modified difference principle, which ranks distributions by the level of the worst-off position and, where that level is the same, by the total. Many actual markets satisfy these principles, and some do not. The thesis has two parts, a criterion and a verdict, and we argue for both. The criterion requires an account of what it is for a market's outcome to satisfy principles that apply to the basic structure as a whole, and Section~3 supplies one. The verdict requires an assessment of actual markets, which we give in outline in Sections~6 and~7.

\section{A short history of the debate}

The modern debate begins with \citet{Smith1776}, who argued that individuals pursuing their own gain in a system of natural liberty are led, as if by an invisible hand, to promote the public interest. Smith did not claim that the resulting distribution is just; he was concerned with the wealth of nations, and he was alert to the ways in which merchants and masters combine against the public and against their workers \citep[I.viii, I.x]{Smith1776}. His successors turned the claim about efficiency into a claim about justice, and each such attempt has been defeated by a case.

\paragraph{Desert by contribution.} The first proposal was that market outcomes are just because they reward each person in proportion to her contribution to production, as measured by the price others will pay for it. The defeating case is the person whose contribution has a high price for reasons unrelated to anything she did: the owner of land that happens to lie over oil, or the worker whose skill is suddenly in demand. Contribution as measured by price depends on scarcity and on the contributions of others, and much of it is luck. POL-19 argues that merit is partly a product of luck and that luck is not deserved; market rewards, which follow merit as the market values it, are therefore not deserved either.

\paragraph{Exploitation.} \citet{Marx1867} argued that capitalist markets are unjust because the worker, who must sell her labour power to live, produces more value than she receives in wages, the difference being appropriated by the owner of capital. The defeating case is a market economy in which the worst off are better off than under any feasible alternative, including every feasible socialist alternative. If in such an economy the surplus is in part reinvested and in part taxed to support those who cannot work, the complaint that the worker produces more than she receives is no longer a complaint that she is treated unjustly, since every alternative that gave her the full product would leave her and others worse off. The exploitation charge has force where it does not depend on the labour theory of value: where workers have no reasonable alternative to accepting terms that a feasible arrangement would improve. That is an outcome complaint, and the criterion of Section~3 accommodates it.

\paragraph{Entitlement.} \citet[ch.~7]{Nozick1974} argued that any patterned principle of justice will be upset by liberty. Begin with a distribution just by any pattern you like; let a million people each pay twenty-five cents to watch Wilt Chamberlain play; the new distribution is unequal, but it arose by voluntary transfers from a just one, so how could it be unjust? The defeating case for the entitlement theory is the starting point. Actual holdings have arisen through conquest, fraud and expropriation, and the principle of rectification that Nozick acknowledged requires information that no one has. The entitlement theory can deliver verdicts about actual markets only if it is supplemented by a principle for the case in which history cannot be recovered, and Nozick himself suggested that such a principle might take the form of the difference principle \citep[ch.~7]{Nozick1974}. We return to Wilt Chamberlain in Section~6.

\paragraph{The mirage of social justice.} \citet{Hayek1945} argued that the price system transmits dispersed knowledge that no central authority could gather, and \citet{Hayek1976} that, since no one designs market outcomes, no one can be unjust in producing them. The defeating case is the law. Markets operate within rules of property, contract and tort, and these rules are chosen, maintained and enforced. Different rules produce different distributions, and the choice among rules is made by agents who can foresee, in general terms, what their choices will produce. That the particular result for a particular person is not intended does not make the choice of rules immune from assessment, any more than the unintended deaths caused by a choice of traffic regulations make that choice immune.

\paragraph{The basic structure.} \citet[\S\S 14, 42--43]{Rawls1971} placed markets within the basic structure: they are part of the system of institutions whose rules distribute the benefits and burdens of social cooperation. Rawls held that the principles of justice apply to the rules, not to particular transactions, and that a market within a just basic structure yields a case of pure background procedural justice: whatever distribution results from the operation of just rules is just \citep{Rawls1977}. He proposed that the institutions of a just society be divided into branches, an allocation branch that keeps markets workably competitive, a stabilization branch that maintains employment, a transfer branch that guarantees a social minimum and a distribution branch that adjusts the background through taxation \citep[\S 43]{Rawls1971}. He later argued that welfare-state capitalism fails his principles and that only a property-owning democracy or liberal socialism could satisfy them \citep[\S\S 41--42]{Rawls2001}, drawing on \citet{Meade1964}. The view defended here is a development of Rawls's, with two differences: we use the modified difference principle of POL-01 rather than his, and we give an explicit criterion for the justice of a single market, which Rawls did not.

\subsection{Recent literature}

The last two decades have produced three lines of work on the justice of markets. The first concerns their moral limits. \citet{Anderson1993} argued that some goods are degraded when they are treated as commodities, \citet{Sandel2012} that markets crowd out non-market norms, and \citet{Satz2010} that certain markets are \emph{noxious}: they produce extreme harms for individuals or for society, or they rest on the extreme vulnerability or weak agency of participants. \citet{BrennanJaworski2016} replied that, if one may give something away, one may sell it, and that the objections to particular markets concern the manner in which they are conducted, not the fact of sale. The second line concerns the place of economic liberty. \citet{Tomasi2012} argued that the economic liberties of ownership and contract should count among the basic liberties, and that a \enquote{market democracy} can satisfy a version of the difference principle; \citet{Freeman2011} replied that the liberal tradition has always treated these liberties as subordinate to distributive principles. The third concerns institutions. \citet{ONeillWilliamson2012} collected work on property-owning democracy, and \citet{Herzog2013} reconstructed the history of the idea of the market in liberal thought, showing that its classical defenders made its justice depend on its institutional setting.

Two further contributions frame the question we ask. \citet{MurphyNagel2002} argued that pre-tax holdings have no moral standing of their own, since property is defined by the legal system, including the tax system: there is no distribution of \enquote{what people would have had} independent of the rules. And \citet{Mirrlees1971} showed that, when abilities are private information, redistribution through taxation faces incentive constraints, so that the feasible set of distributions is smaller than the set obtainable by costless lump-sum transfers. Both results are used below.

\section{Framework}

\begin{definition}\label{d:market}
A \emph{market} is an institution in which goods or services of some kind are allocated through voluntary exchange at prices determined by the exchanges. A market is individuated by the kind of thing exchanged and by its \emph{rules}: the rules of property, contract and regulation that define what may be exchanged, by whom and on what terms.
\end{definition}

\begin{definition}\label{d:regime}
A \emph{regime} is a basic structure: a set of markets together with the background institutions, including the constitution, the tax and transfer system and the provision of public goods. A regime determines an \emph{outcome}, the distribution of primary goods over representative positions, over the course of lives.
\end{definition}

The rules of a market, in Definition~\ref{d:market}, exclude general taxation and transfers, which belong to the background. This follows Rawls's division of branches: the rules of a market are those of the allocation branch as they apply to it, while the transfer and distribution branches are background. The division matters, because otherwise every market could be \enquote{reformed} by adding a tax on its transactions to fund transfers, and the criterion below would collapse into a criterion for the regime.

\begin{principle}[Equal liberty]\label{p:liberty}
Each person has an equal right to the most extensive scheme of equal basic liberties compatible with a similar scheme for all.
\end{principle}

\begin{principle}[Modified difference principle]\label{p:mdp}
Of two outcomes, the better is the one with the higher level for the worst-off position or, if that level is the same, with the higher total.
\end{principle}

These are the principles of POL-01, where Principle~\ref{p:mdp} is Definition~5.1 of POL-01-a. A set of regimes is \emph{feasible} if each can be established given the general facts, including the incentive constraints of \citet{Mirrlees1971}.

\begin{definition}\label{d:justregime}
A regime is \emph{just} if it respects Principle~\ref{p:liberty} and no feasible regime that respects Principle~\ref{p:liberty} has an outcome better by Principle~\ref{p:mdp}.
\end{definition}

\begin{definition}\label{d:justmarket}
A market $m$ in a regime $R$ is \emph{just} if (i) its rules respect Principle~\ref{p:liberty}, and (ii) no feasible regime that differs from $R$ only in the rules of $m$, and respects Principle~\ref{p:liberty}, has an outcome better by Principle~\ref{p:mdp}.
\end{definition}

Definition~\ref{d:justmarket} is the precise form of the thesis that a market is just when its outcome satisfies the principles of POL-01. A market's outcome is the outcome of the regime as determined by the market's rules given the rest of the regime; it satisfies the principles if those rules respect the basic liberties and no alternative choice of them would yield a better outcome. The \enquote{alternatives} include abolishing the market, restricting it, replacing it with public provision and regulating it differently.

\begin{assumption}[Markets are part of the basic structure]\label{a:basic}
The rules of a market are rules of the basic structure: they are coercively enforced, they are chosen and could be otherwise, and their effects on people's life prospects are profound and present from the start.
\end{assumption}

\begin{assumption}[Background procedural justice]\label{a:background}
Within a regime whose institutions satisfy the principles of justice, the distribution that results from people's choices is just, whatever it is, and particular transactions are not separately subject to the principles.
\end{assumption}

Assumption~\ref{a:basic} is Rawls's, and the reply to Hayek in Section~2 is its defence: the rules of property and contract are chosen, and their choice can be assessed by what they produce. Assumption~\ref{a:background} is the reason why the principles are applied to rules and not to transactions; it is defended in \citet{Rawls1977} and is the answer to Wilt Chamberlain (Section~6).

\section{The criterion}

\begin{theorem}[The outcome criterion]\label{t:criterion}
Whether a market is just depends only on whether its rules respect the basic liberties and on the outcomes that the feasible choices of its rules produce. In particular, the voluntariness of the exchanges in a market is neither necessary nor sufficient for its justice.
\end{theorem}

\begin{proof}
By Assumption~\ref{a:basic}, the rules of a market are part of the basic structure, and by POL-01 the justice of the basic structure is fixed by Principles~\ref{p:liberty} and~\ref{p:mdp}; Definition~\ref{d:justmarket} applies them to the choice of a market's rules. By Assumption~\ref{a:background}, nothing about particular transactions adds a further condition. So the justice of a market depends only on the two factors stated.

That voluntariness is not sufficient: consider two regimes whose markets have the same rules, so that every exchange in each is equally voluntary, and which differ in the background distribution of holdings because one of them has a transfer system and the other does not. Let the first produce an outcome with worst-off level~$2$ and total~$20$, and the second an outcome with worst-off level~$4$ and total~$19$. By Principle~\ref{p:mdp}, the second is better, and the first regime is unjust, though its exchanges are as voluntary as those of the second. Holding the background fixed, a market whose exchanges are voluntary can also be unjust, if a feasible change in its rules, such as a minimum standard for the terms of employment, raises the worst-off level.

That voluntariness is not necessary: a market may include rules that restrict which exchanges may be made, such as licensing requirements, prohibitions of certain contracts or price regulation, and so limit voluntary exchange. If those rules respect the basic liberties and no feasible alternative to them yields a better outcome, the market is just by Definition~\ref{d:justmarket}.
\end{proof}

Theorem~\ref{t:criterion} disposes of two of the three answers of Section~1. Against the entitlement view, voluntariness does not make an outcome just; against the view that markets are unjust in themselves, the market form is not a ground of injustice, since whether a market is just depends on what its rules produce, and in many cases the market form produces outcomes that no feasible alternative improves. The third answer, that justice does not apply to markets, is denied by Assumption~\ref{a:basic}.

\section{The structure of market justice}

The criterion has several structural consequences.

\begin{proposition}\label{p:global}
If a regime is just, every market in it is just.
\end{proposition}

\begin{proof}
If a market $m$ in a just regime $R$ were unjust, then either its rules would violate Principle~\ref{p:liberty}, in which case $R$ would not respect it, or some feasible regime differing from $R$ only in the rules of $m$ would have a better outcome, in which case $R$ would not be just by Definition~\ref{d:justregime}.
\end{proof}

\begin{proposition}\label{p:local}
The converse fails: every market in a regime may be just although the regime is unjust.
\end{proposition}

\begin{proof}
The injustice of a regime may lie in its background institutions. Suppose that the markets of a regime have the best rules available given its tax and transfer system, but that a different tax and transfer system would raise the worst-off level. Then no change in the rules of any single market yields a better outcome, so each market is just, while the regime is not. The same holds if the only improvements require changing the rules of two markets at once.
\end{proof}

Proposition~\ref{p:local} is what makes the second part of our thesis compatible with the evident injustice of many actual societies. A market can be just in an unjust society, because the injustice is located elsewhere: in the tax system, in the provision of education, in the absence of a social minimum. Conversely, Proposition~\ref{p:global} ensures that the assessment of single markets is not idle: in a just society every market passes it.

\begin{proposition}\label{p:pareto}
If a regime is just, its outcome is Pareto efficient among the outcomes of feasible regimes that respect Principle~\ref{p:liberty}: no such regime makes some position better off and none worse off.
\end{proposition}

\begin{proof}
Suppose that a feasible regime makes some position better off and none worse off. Its worst-off level is at least as high, since no position is worse off, and its total is strictly higher. By Principle~\ref{p:mdp} its outcome is better, either because the worst-off level is higher or, if it is the same, because the total is. So the original regime is not just.
\end{proof}

\begin{proposition}\label{p:notconverse}
Pareto efficiency does not imply justice.
\end{proposition}

\begin{proof}
Let two feasible outcomes have worst-off levels $1$ and $3$ and totals $10$ and $6$. If the first gives a lower level to the worst-off position and a higher one to another, neither outcome Pareto dominates the other, and both may be efficient. By Principle~\ref{p:mdp}, the second is better, so the first is not just.
\end{proof}

Propositions~\ref{p:pareto} and~\ref{p:notconverse} locate the classical welfare theorems in the theory of justice. The first welfare theorem says that competitive equilibria are Pareto efficient under certain conditions \citep{Arrow1951,Debreu1959}. By Proposition~\ref{p:notconverse}, this is not a theorem about justice. The second welfare theorem says that any efficient allocation can be reached as a competitive equilibrium after suitable lump-sum transfers. If the just outcome is efficient, as Proposition~\ref{p:pareto} says it is, the second theorem shows that markets are, in principle, a means of reaching it. The qualification \enquote{in principle} matters: lump-sum transfers based on ability are not feasible when ability is private, and the feasible set is constrained accordingly \citep{Mirrlees1971}. The relation between entitlement theory and lump-sum transfers is the subject of POL-01-c.

\begin{proposition}\label{p:liberties}
A market whose objects are basic liberties, or whose rules permit a person to alienate them, is unjust in every regime.
\end{proposition}

\begin{proof}
Such a market violates Principle~\ref{p:liberty}, and condition (i) of Definition~\ref{d:justmarket} fails whatever the rest of the regime is.
\end{proof}

Proposition~\ref{p:liberties} is the place where the priority of liberty enters the assessment of markets. Markets in votes, in political office, in the services of judges and in persons fail it, as do labour contracts that alienate the right to leave. They fail it whatever their outcomes for the worst off.

\paragraph{Sufficiency and equality.} The criterion does not condemn inequality as such. The modified difference principle permits an inequality that raises the worst-off level, or that leaves it unchanged and increases the total, and it never prefers an outcome produced by levelling down. POL-08 argues that equality is not valuable in itself and that what matters is that everyone has enough. The two views agree on markets: the floor that POL-01 presupposes, a level adequate for a decent life and for the exercise of the basic liberties (Assumption~3.3 of POL-01-a), is the sufficiency threshold of POL-08, and the modified principle is prioritarian at the bottom, not egalitarian. A market that produces great inequality while everyone has enough is not, for that reason, unjust; it is unjust if a feasible change in its rules would raise the worst-off level, or keep it and raise the total.

\section{Cases}

\paragraph{Wilt Chamberlain.} Nozick's case is a challenge to any principle that judges outcomes. The challenge assumes that the principle is applied to a snapshot of holdings and that maintaining it requires interfering with each transaction. Neither assumption holds for Definition~\ref{d:justmarket}. The principles apply to the rules of the regime, which include the market for basketball tickets and a tax system, and the outcome they assess is the distribution over lives that results from people's choices under those rules. If the regime is just, the distribution after the season, including Chamberlain's income net of the taxes that the rules impose, is just by Assumption~\ref{a:background}. No transaction is interfered with; the tax is a known feature of the rules under which Chamberlain contracted. Whether the tax is just is determined by Principle~\ref{p:mdp}: it is just if no feasible alternative rule raises the worst-off level, or keeps it and raises the total. Nozick's question \enquote{how could the new distribution be unjust?} has an answer: it could be unjust only if the rules under which it arose are.

\paragraph{Sweatshop labour.} Workers in export factories in poor countries accept wages and conditions that would be unlawful in rich ones. \citet{Zwolinski2007} argued that, since the workers choose these jobs over their alternatives, prohibiting the jobs or imposing standards that close them makes the workers worse off. Our criterion asks two questions. The first is whether the rules respect the basic liberties: a factory that confines its workers, or a contract that cannot be left, fails Proposition~\ref{p:liberties}. The second is whether a feasible change in the rules, such as enforceable safety standards or a right to organize, would raise the level of the worst off. If it would, the market is unjust, whatever the workers' choices; if every feasible standard would close the factories and leave the workers worse off, it is not. The answer depends on facts that vary between countries and industries, and on whether the relevant regime is national or global; POL-13 argues that duties of beneficence apply globally but require international institutions for their fulfilment, and the regime against which these markets are assessed includes those institutions.

\paragraph{Votes.} A market in votes may raise everyone's income, if the poor sell their votes for more than the policies they would have voted for are worth to them. It is nonetheless unjust in every regime, by Proposition~\ref{p:liberties}: the equal political liberties are basic, and their fair value is destroyed when they can be bought.

\paragraph{Organs.} The kidney market is the standard case in the literature on noxious markets \citep{Satz2010}. On the present criterion its justice is conditional. The right to bodily integrity is a basic liberty, and a market that permits its alienation under duress fails Proposition~\ref{p:liberties}; but sale is not alienation of the right, any more than paid employment in a dangerous occupation is. The second question is whether a regulated market, with a single public purchaser, fixed compensation and medical safeguards, would produce a better outcome for the worst off than the feasible alternatives, which include the present system of donation and the illegal trade it sustains. That is an empirical question, and the criterion tells us what would settle it.

\paragraph{Famine.} \citet{Sen1981} showed that famines have occurred without any decline in the availability of food, when the market value of what the poor had to sell, their labour or their animals, collapsed relative to the price of food. A food market that operates in this way pushes some below the floor, and feasible changes in its rules, such as public employment that maintains purchasing power or the release of public stocks, would raise the worst-off level. Such a market is unjust by Definition~\ref{d:justmarket}, though every exchange in it is voluntary.

\paragraph{Winner-take-all labour markets.} In some labour markets small differences in performance yield very large differences in reward \citep{Frank2016}. These rewards are not deserved (POL-19), but that does not make the market unjust: by Theorem~\ref{t:criterion}, the question is whether a feasible change in the rules would improve the outcome by Principle~\ref{p:mdp}. Where the large rewards are incentive payments that elicit effort or skill whose products benefit the worst off, the market may be just, and the inequality is a legitimate expectation, not a desert. Where the rewards are rents, which a change in the rules of competition or of corporate governance could reduce without lowering anyone's level, the market is unjust, because the change would raise the total without lowering the worst-off level, or would free resources for the background institutions to raise it.

\section{Many, but not all}

The verdict that many actual markets are just rests on Proposition~\ref{p:local} and on two general facts. The first is the informational argument of \citet{Hayek1945}: in markets for ordinary goods and services, prices coordinate the use of dispersed knowledge, and the feasible alternatives to market allocation, central allocation or rationing, perform worse, for everyone and in particular for the worst off. The historical experience of planned economies in consumer goods supports this. The second is that the injustices of actual societies are, in the main, located in their background institutions: in insufficient transfers, in unequal access to education and health care, in unequal political influence. By Proposition~\ref{p:local}, these do not make the markets for bread, clothing, books, furniture or restaurant meals unjust, since no change in the rules of those markets would remedy them. For such markets, the rules typical of developed economies, with contract law, consumer protection and competition policy, are plausibly the best feasible rules given the background, and they are just in the sense of Definition~\ref{d:justmarket}.

The verdict that not all actual markets are just rests on the cases of Section~6 and on the structure of the criterion. Some actual markets violate the basic liberties, as forced labour and the sale of political influence do. Others leave some below the floor where a feasible change in their rules would prevent it, as the food markets of famines did. Others yield rents that a feasible change in their rules would reduce without harm to anyone. The criterion identifies these failures without condemning the market form, and it identifies the remedy: a change in the market's rules where the defect is in them, a change in the background where it is not. The companion manuscript POL-12-b assesses particular actual markets.

\section{Objections}

\begin{objection}[Taxation and liberty]
The criterion permits taxing the proceeds of voluntary exchange to raise the worst-off level. Taxation of earnings is on a par with forced labour, since it takes the product of some hours of a person's work \citep[ch.~7]{Nozick1974}.
\end{objection}

\begin{reply}
The objection assumes that a person's pre-tax income is hers independently of the rules, so that taxation takes what she already owns. But property is defined by the legal system, and the tax system is part of it \citep{MurphyNagel2002}; there is no pre-tax entitlement against which taxation can be measured. Taxation does not compel anyone to work, and a person who chooses not to work is not taxed for the hours she does not work. POL-06 argues that property is not theft; the converse thesis, that taxation is theft, fails for a parallel reason, since it presupposes an entitlement that only the rules, taxes included, can define.
\end{reply}

\begin{objection}[Spontaneous order]
No one designs the distribution that a market produces. Calling it unjust is like calling the weather unjust \citep{Hayek1976}.
\end{objection}

\begin{reply}
The weather is not produced by rules that anyone could change. The distribution that a market produces depends on the rules of property, contract and taxation, which are chosen, and Definition~\ref{d:justmarket} assesses the choice. That no one intends the particular outcome for a particular person is compatible with the choice of rules being foreseeably better or worse for the worst off, and that foreseeability is what makes the choice assessable.
\end{reply}

\begin{objection}[Incentives and ethos]
A just society would not need incentive inequalities. The talented demand higher rewards as a condition of working productively, and a society whose members accepted the difference principle would not make the demand \citep{Cohen2008}. Markets that pay incentives are unjust because they accommodate a selfish ethos.
\end{objection}

\begin{reply}
The principles of POL-01 apply to the basic structure, and the feasible set is fixed by the facts about motivation that obtain, including the motivations that markets channel. Cohen's argument shows that a society of persons with a different ethos could do better; it does not show that the rules are unjust given the persons who live under them. A market that pays incentives is just if no feasible change in its rules, with motivations as they are, raises the worst-off level. Whether individuals have reasons of justice, in their personal choices, to accept less than the market offers is a question about personal ethics that our criterion leaves open.
\end{reply}

\begin{objection}[The moral limits of markets]
Some goods are corrupted by sale: friendship, sex, citizenship, children. The criterion, which considers only liberty and outcomes, cannot explain why markets in these goods are wrong \citep{Anderson1993,Sandel2012}.
\end{objection}

\begin{reply}
Several of the goods in question are covered by Proposition~\ref{p:liberties}: citizenship and political influence are basic liberties, and the sale of children treats persons as objects of property. Of the remaining cases, some are covered by the outcome condition, since markets that exploit extreme vulnerability tend to leave the vulnerable worse off than feasible alternatives would, which is the core of Satz's account of noxious markets \citep{Satz2010}. What remains is the claim that some exchanges express inappropriate attitudes toward their objects. That claim may be true, but it concerns the value of the goods and the virtue of the parties, not justice, and the thesis that markets are just when their outcomes satisfy the principles is a thesis about justice.
\end{reply}

\begin{objection}[The verdict is unverifiable]
Whether a given market's rules are the best feasible rules depends on counterfactual facts that no one can establish, so the claim that many actual markets are just is not supported.
\end{objection}

\begin{reply}
The criterion makes the verdict depend on counterfactual facts, as every outcome-based criterion must, but the facts are not inaccessible. For many markets, the alternatives have been tried: price controls, rationing, central allocation, prohibition. The evidence from these trials, together with the informational argument, supports the verdict for ordinary markets. The verdict is defeasible, and the criterion says what would defeat it in any case: evidence that a feasible change in the rules would raise the level of the worst off.
\end{reply}

\section{What the result does not establish}

The argument does not establish which of the regimes Rawls considered, property-owning democracy, liberal socialism or a reformed welfare-state capitalism, is the just regime for a given society; the institutional consequences of POL-01 are derived in its second volume. It does not establish the verdict for any particular market in detail, which requires empirical work. Nor does it extend the criterion to global markets, whose regime is not a single basic structure. Finally, it does not show that market rewards are deserved; they are legitimate expectations under just rules, and POL-19 argues that the merit they track is partly a product of luck.

\section{Conclusion}

A market is a component of the basic structure, defined by rules that are chosen and enforced. It is just when its rules respect the basic liberties and no feasible choice of them would give the worst off more, or give them as much and everyone else more. The voluntariness of exchange does not make a market just, and the market form does not make it unjust; justice in the whole regime implies justice in each market, and efficiency is necessary but not sufficient for it. Applied to actual markets, the criterion vindicates many, such as the ordinary markets in goods and services of developed economies, whose rules no feasible change would improve for the worst off, and condemns some, such as markets in political influence, the food markets of famines and labour markets that yield removable rents.

\begin{thebibliography}{99}\small
\bibitem[Anderson(1993)]{Anderson1993} Anderson, E. (1993). \emph{Value in Ethics and Economics}. Harvard University Press.
\bibitem[Arrow(1951)]{Arrow1951} Arrow, K.~J. (1951). An extension of the basic theorems of classical welfare economics. In J. Neyman (ed.), \emph{Proceedings of the Second Berkeley Symposium on Mathematical Statistics and Probability}. University of California Press.
\bibitem[Brennan and Jaworski(2016)]{BrennanJaworski2016} Brennan, J. and Jaworski, P.~M. (2016). \emph{Markets without Limits: Moral Virtues and Commercial Interests}. New York: Routledge.
\bibitem[Cohen(2008)]{Cohen2008} Cohen, G.~A. (2008). \emph{Rescuing Justice and Equality}. Harvard University Press.
\bibitem[Cohen(2009)]{Cohen2009} Cohen, G.~A. (2009). \emph{Why Not Socialism?} Princeton University Press.
\bibitem[Debreu(1959)]{Debreu1959} Debreu, G. (1959). \emph{Theory of Value: An Axiomatic Analysis of Economic Equilibrium}. New York: Wiley.
\bibitem[Frank(2016)]{Frank2016} Frank, R.~H. (2016). \emph{Success and Luck: Good Fortune and the Myth of Meritocracy}. Princeton University Press.
\bibitem[Freeman(2011)]{Freeman2011} Freeman, S. (2011). Capitalism in the classical and high liberal traditions. \emph{Social Philosophy and Policy}, 28(2), 19--55.
\bibitem[Hayek(1945)]{Hayek1945} Hayek, F.~A. (1945). The use of knowledge in society. \emph{American Economic Review}, 35(4), 519--530.
\bibitem[Hayek(1976)]{Hayek1976} Hayek, F.~A. (1976). \emph{Law, Legislation and Liberty, Volume 2: The Mirage of Social Justice}. University of Chicago Press.
\bibitem[Herzog(2013)]{Herzog2013} Herzog, L. (2013). \emph{Inventing the Market: Smith, Hegel, and Political Theory}. Oxford University Press.
\bibitem[Marx(1867)]{Marx1867} Marx, K. (1867). \emph{Das Kapital: Kritik der politischen \"Okonomie}, vol.~1. Hamburg: Otto Meissner.
\bibitem[Meade(1964)]{Meade1964} Meade, J.~E. (1964). \emph{Efficiency, Equality and the Ownership of Property}. London: Allen and Unwin.
\bibitem[Mirrlees(1971)]{Mirrlees1971} Mirrlees, J.~A. (1971). An exploration in the theory of optimum income taxation. \emph{Review of Economic Studies}, 38(2), 175--208.
\bibitem[Murphy and Nagel(2002)]{MurphyNagel2002} Murphy, L. and Nagel, T. (2002). \emph{The Myth of Ownership: Taxes and Justice}. Oxford University Press.
\bibitem[Nozick(1974)]{Nozick1974} Nozick, R. (1974). \emph{Anarchy, State, and Utopia}. New York: Basic Books.
\bibitem[O'Neill and Williamson(2012)]{ONeillWilliamson2012} O'Neill, M. and Williamson, T. (eds.) (2012). \emph{Property-Owning Democracy: Rawls and Beyond}. Chichester: Wiley-Blackwell.
\bibitem[Rawls(1971)]{Rawls1971} Rawls, J. (1971). \emph{A Theory of Justice}. Harvard University Press.
\bibitem[Rawls(1977)]{Rawls1977} Rawls, J. (1977). The basic structure as subject. \emph{American Philosophical Quarterly}, 14(2), 159--165.
\bibitem[Rawls(2001)]{Rawls2001} Rawls, J. (2001). \emph{Justice as Fairness: A Restatement}. Ed. E. Kelly. Harvard University Press.
\bibitem[Sandel(2012)]{Sandel2012} Sandel, M.~J. (2012). \emph{What Money Can't Buy: The Moral Limits of Markets}. New York: Farrar, Straus and Giroux.
\bibitem[Satz(2010)]{Satz2010} Satz, D. (2010). \emph{Why Some Things Should Not Be for Sale: The Moral Limits of Markets}. Oxford University Press.
\bibitem[Sen(1981)]{Sen1981} Sen, A. (1981). \emph{Poverty and Famines: An Essay on Entitlement and Deprivation}. Oxford: Clarendon Press.
\bibitem[Smith(1776)]{Smith1776} Smith, A. (1776). \emph{An Inquiry into the Nature and Causes of the Wealth of Nations}. London: Strahan and Cadell. Cited by book and chapter.
\bibitem[Tomasi(2012)]{Tomasi2012} Tomasi, J. (2012). \emph{Free Market Fairness}. Princeton University Press.
\bibitem[Zwolinski(2007)]{Zwolinski2007} Zwolinski, M. (2007). Sweatshops, choice, and exploitation. \emph{Business Ethics Quarterly}, 17(4), 689--727.
\end{thebibliography}

\end{document}
